\sheet{4}{Sampling}{Chapter~4 (Sampling and Quantisation)}

\begin{goals}
\begin{itemize}
  \item Derive and apply the quantisation-noise model and the rule
    $\mathrm{SNR}\approx 6.02N+1.76$\,dB; size an ADC for the knock sensor.
  \item Configure timer-driven sampling on a microcontroller and verify the
    sampling rate and the conversion time with the logic analyzer.
  \item Quantify the effect of sampling-time jitter and compare it with
    quantisation noise.
  \item Resample a time-sampled engine signal onto a crank-angle grid using
    the edges of the 60-2 crank signal, and explain when the time domain and
    when the angle domain is the natural one.
\end{itemize}
\end{goals}

\section*{Background}

An ADC with $N$ bits and full-scale range $V_\mathrm{FS}$ has the step size
(1 LSB) $\Delta = V_\mathrm{FS}/2^N$. If the signal is ``busy'' (it moves
over many steps between samples), the quantisation error
$e=x_q-x$ is well modelled as white noise, uniformly distributed in
$[-\Delta/2,\Delta/2)$. Sampling instants that deviate randomly from the
ideal grid, $t_n = nT_s+\delta_n$ with $\operatorname{std}(\delta_n)=\sigma_j$
(\emph{jitter}), produce an additional error $e_n\approx\dot x(nT_s)\,\delta_n$
that grows with the signal frequency.

An ECU samples two kinds of signals: \emph{time-based} (fixed $T_s$, e.g.\ the
knock sensor with \SIrange{20}{50}{kHz}) and \emph{angle-based} (one value per
fixed crank-angle increment, triggered by the crank wheel). An angle grid can
also be computed afterwards from time samples: the crank edges provide
pairs (time, angle), between which the angle is interpolated.

%-------------------------------------------------------------------------
\begin{exercise}{Quantisation and jitter budget}{\Paper}
\begin{tasks}
  \item Assume a uniformly distributed quantisation error. Derive its
    variance and show that a full-scale sine quantised with $N$ bits has
    $\mathrm{SNR}=6.02N+1.76$\,dB. By how much does the SNR drop if the sine
    uses only $1/10$ of the full-scale amplitude?
  \item Compute 1\,LSB and the rms quantisation noise for the Arduino Uno
    (10 bit, \SIrange{0}{5}{V}) and the ESP32 (12 bit, \SIrange{0}{3.3}{V}).
  \item The conditioned knock-sensor signal spans \SIrange{0}{5}{V} (offset
    \SI{2.5}{V}) for the heaviest knock. The weakest knock that still has to be
    measured has an amplitude of \SI{10}{mV}; its quantisation SNR must be at
    least \SI{30}{dB}. Which dynamic range (in dB) is required, and how many
    bits does the ADC need?
  \item The knock detector only uses a band of width $B=\SI{1}{kHz}$ around
    the knock frequency, the ADC samples at $f_s=\SI{50}{kHz}$. Assuming white
    quantisation noise, by how many dB does a band-pass filter improve the SNR
    of (c)? How many ADC bits does this ``save''?
  \item Show that jitter limits the SNR of a sine of frequency $f$ to
    $\mathrm{SNR}_j=-20\log_{10}(2\pi f\sigma_j)$. (i) Which jitter is allowed
    for a \SI{6.5}{kHz} knock signal if jitter shall not degrade a 12-bit ADC?
    (ii) A sampling loop that polls \texttt{micros()} (resolution \SI{4}{\micro
    s} on the Uno) produces a timing error uniformly distributed over
    \SI{4}{\micro s}. Compute $\sigma_j$, $\mathrm{SNR}_j$ at \SI{6.5}{kHz}
    and the \emph{effective number of bits}
    $\mathrm{ENOB}=(\mathrm{SNR}-1.76)/6.02$.
\end{tasks}
\end{exercise}

\begin{solution}
a) For $e\sim\mathcal{U}[-\Delta/2,\Delta/2)$:
$\sigma_e^2=\frac1\Delta\int_{-\Delta/2}^{\Delta/2}e^2\dd e=\frac{\Delta^2}{12}$.
A full-scale sine has amplitude $A=2^{N-1}\Delta$ and power
$P_s=A^2/2=2^{2N}\Delta^2/8$. Hence
\[
\mathrm{SNR}=\frac{P_s}{\sigma_e^2}=\frac{3}{2}\,2^{2N}
\;\Rightarrow\;
10\log_{10}\mathrm{SNR}=10\log_{10}1.5+20N\log_{10}2=1.76+6.02N\ \text{dB}.
\]
The noise power does not depend on the signal, so an amplitude of $1/10$ of
full scale lowers the SNR by $20\log_{10}10=\SI{20}{dB}$ (equivalent to losing
$3.3$ bits). Every unused part of the input range costs resolution.

b) Uno: $\Delta=5/1024\,\text{V}=\SI{4.88}{mV}$,
$\sigma_e=\Delta/\sqrt{12}=\SI{1.41}{mV}$.
ESP32: $\Delta=3.3/4096\,\text{V}=\SI{0.806}{mV}$, $\sigma_e=\SI{0.233}{mV}$.

c) Dynamic range: $20\log_{10}(\SI{2.5}{V}/\SI{10}{mV})=\SI{48}{dB}$.
A sine of amplitude $a$ has $\mathrm{SNR}=(a^2/2)/(\Delta^2/12)=6a^2/\Delta^2$.
$\mathrm{SNR}\ge 10^3$ requires $\Delta\le a\sqrt{6\cdot10^{-3}}=\SI{0.775}{mV}$,
i.e.\ $2^N\ge \SI{5}{V}/\SI{0.775}{mV}=6455$, so $N=13$ bits (a 12-bit ADC
gives $\Delta=\SI{1.22}{mV}$ and only \SI{26}{dB}). Equivalently:
$48+30=\SI{78}{dB}$ of ``full-scale SNR'' $\Rightarrow N\ge(78-1.76)/6.02=12.7$.

d) White quantisation noise is spread evenly over $[0,f_s/2]$; a filter of
bandwidth $B$ keeps the fraction $2B/f_s=1/25$. The SNR improves by
$10\log_{10}(f_s/2B)=\SI{14}{dB}$, i.e.\ $14/6.02\approx2.3$ bits: with the
band-pass a 12-bit ADC reaches $26+14=\SI{40}{dB}$ and even 11 bits
($\approx\SI{34}{dB}$) would suffice. (Oversampling and filtering trade
sampling rate for resolution: \SI{3}{dB}, i.e.\ half a bit, per doubling of
$f_s$.)

e) For $x=A\sin\omega t$: $e_n\approx A\omega\cos(\omega nT_s)\,\delta_n$ with
independent $\delta_n$, so $\overline{e^2}=A^2\omega^2\sigma_j^2/2$ and
$\mathrm{SNR}_j=(A^2/2)/\overline{e^2}=1/(\omega\sigma_j)^2$, i.e.\
$\mathrm{SNR}_j=-20\log_{10}(2\pi f\sigma_j)$.
(i) 12 bit: $\mathrm{SNR}_q=\SI{74.0}{dB}$; $2\pi f\sigma_j\le10^{-74/20}=2.0\cdot10^{-4}$
gives $\sigma_j\le\SI{4.9}{ns}$ at \SI{6.5}{kHz} -- only hardware-triggered
conversions achieve this.
(ii) Uniform over \SI{4}{\micro s}: $\sigma_j=\SI{4}{\micro s}/\sqrt{12}=\SI{1.15}{\micro s}$,
$\mathrm{SNR}_j=-20\log_{10}(2\pi\cdot6500\cdot1.15\cdot10^{-6})=\SI{26.5}{dB}$,
$\mathrm{ENOB}=4.1$ bits. A software-polled sampling loop turns a 10-bit ADC
into a 4-bit ADC for knock signals!
\end{solution}

%-------------------------------------------------------------------------
\begin{exercise}{ADC in practice: timer-driven sampling}{\Wokwi}
Open the Uno project \codefile{wokwi/sheet04\_adc/} (sketch, diagram and the
\texttt{engine-sim} chip files). Timer1 triggers an interrupt service routine
(ISR) at $f_s=\SI{5}{kHz}$ that samples the \texttt{KNOCK} output of the chip
(A1) into a RAM buffer of 400 samples; the buffer is printed afterwards (the
UART cannot transmit \SI{5}{kHz}\,$\times$ a CSV line in real time). Pin 8
toggles at every sample (logic analyzer D0), pin 9 is HIGH while the ISR runs
(D1), the chip's \texttt{TDC} output (pin 2, INT0) marks combustion TDCs.
\begin{tasks}
  \item Timer1 runs in CTC mode with prescaler 8. Compute the compare value
    \texttt{OCR1A} for \SI{5}{kHz} and complete \texttt{timer1\_start()}. What
    is the lowest sampling rate possible with this prescaler?
  \item Complete the ISR. Record about \SI{0.1}{s} with the logic analyzer
    and determine from the VCD file (i) the sampling period (D0) and (ii) the
    time the ISR needs (pulse width on D1). Which sampling rate is the
    theoretical maximum for this structure? Why is the practical limit lower?
  \item Between combustions, \SIrange{140}{175}{\degree} after each TDC, the
    knock signal contains only sensor noise. Complete the computation of the
    index range of this quiet window (at \SI{3000}{rpm}) and run the sketch
    with the noise slider at \SI{50}{mV}, \SI{10}{mV} and \SI{0}{mV}. Compare
    the printed standard deviation (in LSB) with your expectation
    (chip output = $\SI{2.5}{V}+0.6\cdot$ sensor signal). Would the 12-bit ESP32
    ADC improve the knock measurement here?
  \item Turn the potentiometer to an arbitrary position and read the code
    several times. Is the quantisation error of this constant input
    ``noise''? What does this mean for the validity of the model of
    Exercise~4.1?
\end{tasks}
\begin{hint}
Wokwi measures time in \emph{simulated} time; the Uno core is simulated
instruction by instruction, but treat the measured durations as indicative.
\end{hint}
\end{exercise}

\begin{solution}
a) Timer clock $f_\mathrm{T}=\SI{16}{MHz}/8=\SI{2}{MHz}$; the period is
$(\texttt{OCR1A}+1)/f_\mathrm{T}$, so $\texttt{OCR1A}=2\cdot10^6/5000-1=399$.
The 16-bit register allows at most $65536/f_\mathrm{T}=\SI{32.8}{ms}$, i.e.\
$f_{s,\min}=\SI{30.5}{Hz}$ (prescaler 64/256/1024 for slower rates).
\inputcode[firstline=29,lastline=52]{solutions/wokwi/sheet04_adc/sketch.ino}

b) D0 changes level every \SI{200}{\micro s} (square wave with period
\SI{400}{\micro s}). The D1 pulses last about \SI{110}{\micro s}: the
conversion itself takes 13 ADC clock cycles at $\SI{16}{MHz}/128=\SI{125}{kHz}$,
i.e.\ \SI{104}{\micro s}, plus the overhead of \texttt{analogRead()} and the ISR
entry/exit. Theoretical maximum: $1/\SI{112}{\micro s}\approx\SI{8.9}{kHz}$ --
then the CPU does nothing but sampling. In practice other interrupts (Timer0
for \texttt{millis()}, INT0, UART) must fit in between and the main program
needs CPU time, so about \SIrange{5}{7}{kHz} is a sensible limit with
\texttt{analogRead()}. Single edges on D0 may be delayed by a few
\si{\micro s} when the Timer0 ISR is running at the compare match: this is
jitter (it can be removed by disabling Timer0, see Sheet~5).

c) Samples per degree: $f_s/(6\,n)=5000/18000=0.278$, so
$i_0=\lceil140\cdot0.278\rceil=39$, $i_1=\lfloor175\cdot0.278\rfloor=48$
(10 samples per TDC, about 70 per block). Expected standard deviation:
\begin{center}\small
\begin{tabular}{@{}lccc@{}}\toprule
noise slider & sensor noise at ADC & expected std & simulation (3990 samples)\\\midrule
\SI{50}{mV} & \SI{30}{mV} & \SI{6.15}{LSB} & \SI{6.06}{LSB}\\
\SI{10}{mV} & \SI{6}{mV}  & \SI{1.26}{LSB} & \SI{1.25}{LSB}\\
\SI{0}{mV}  & 0 & 0 & \SI{0.46}{LSB}\\\bottomrule
\end{tabular}
\end{center}
(``expected'' $=\sqrt{(0.6\sigma_\mathrm{noise}/\Delta)^2+1/12}$; the
simulation column was obtained with the same engine model sampled exactly like
the sketch. A single block contains only $\approx70$ samples, so the printed
value scatters by about $\pm9\,\%$.) The quantisation noise
($0.29$\,LSB) is negligible against the sensor noise
($\sqrt{6.14^2+0.29^2}=6.15$, i.e.\ \SI{0.01}{dB} worse): the
measurement is limited by the \emph{sensor}, not by the ADC; a 12-bit ADC
would not help. Only for the noise-free signal does quantisation dominate:
the offset \SI{2.5}{V} lies exactly on the transition between codes 511 and
512, so tiny residues of the decaying bursts make the LSB flicker
(std $\approx0.5$\,LSB in the simulation; in Wokwi the result depends on how
the simulator rounds and may be exactly 0).

d) The code is constant, e.g.\ \SI{3.30}{V} gives
$\lfloor 3.30\cdot1024/5\rfloor=675$ every time. The error is a fixed, signal
dependent offset (up to 1 LSB), not a random variable: the noise model of
4.1 requires a signal that changes over many LSB between samples (or added
noise/dither). For slowly varying signals the quantisation error is
correlated and appears as distortion, not as white noise.
\end{solution}

%-------------------------------------------------------------------------
\begin{exercise}{Quantisation noise and jitter}{\JSLinux}
\codefile{jslinux/sheet04/quant.c} generates a sine of
$f_0=\SI{1234.5}{Hz}$ at $f_s=\SI{50}{kHz}$ ($10^5$ samples; $f_0/f_s$ is not
a simple fraction, so the error is not periodic). The full-scale range is
$[-1,1)$.
\begin{tasks}
  \item Implement the function \texttt{quantise()}, a mid-tread quantiser
    with $N$ bits: rounding to the nearest step and clipping to the codes
    $q=-2^{N-1},\dots,2^{N-1}-1$. Implement \texttt{snr\_db()}. Measure the SNR for $N=4\dots16$ at full scale and
    compare with $6.02N+1.76$\,dB. For which $N$ does the formula fail, and
    why?
  \item Repeat with an amplitude of 0.1 (\SI{-20}{dBFS}).
  \item Implement jittered sampling (Gaussian $\delta_n$, no quantisation)
    and measure the SNR for $f=\SI{100}{Hz}\dots\SI{20}{kHz}$ and
    $\sigma_j=0.1,1,\SI{4}{\micro s}$. Compare with
    $-20\log_{10}(2\pi f\sigma_j)$.
  \item Combine a 12-bit quantiser with jitter: compute the ENOB for the
    knock frequency \SI{6.5}{kHz}. At which signal frequency are jitter
    noise and quantisation noise equal for $\sigma_j=\SI{1}{\micro s}$?
\end{tasks}
\end{exercise}

\begin{solution}
\inputcode[firstline=41,lastline=61]{solutions/jslinux/sheet04/quant.c}
Jittered sample: \lstinline|y[n] = A*sin(2*M_PI*f*(t + sj*grand()) + 0.3);|

a/b) Program output (SNR in dB):
\begin{center}\small
\begin{tabular}{@{}rrrrr@{}}\toprule
$N$ & full scale & $6.02N+1.76$ & \SI{-20}{dBFS} & expected\\\midrule
4 & 24.74 & 25.84 & 5.27 & 5.84\\
5 & 31.20 & 31.86 & 10.56 & 11.86\\
6 & 37.47 & 37.88 & 18.80 & 17.88\\
8 & 49.75 & 49.92 & 29.78 & 29.92\\
10 & 61.89 & 61.96 & 42.18 & 41.96\\
12 & 73.97 & 74.00 & 53.97 & 54.00\\
14 & 86.03 & 86.04 & 66.10 & 66.04\\
16 & 98.10 & 98.08 & 78.08 & 78.08\\\bottomrule
\end{tabular}
\end{center}
From 8 bits on, the formula holds within \SI{0.2}{dB}. For $N\le5$ the
measured SNR is up to \SI{1.1}{dB} lower: with few levels the error is no
longer uniform and independent of the signal (it is a deterministic,
signal-dependent staircase error with harmonics), and the amplitude
$1-\Delta/2$ needed to avoid clipping reduces the signal power. The
\SI{-20}{dBFS} column is \SI{20}{dB} lower (within $\pm\SI{0.4}{dB}$) for $N\ge7$; for small $N$ the
reduced sine spans only a few codes and the deviations are larger
($\pm\SI{1.3}{dB}$).

c) Measured vs.\ theory (dB):
\begin{center}\small
\begin{tabular}{@{}r|rr|rr|rr@{}}\toprule
$f$ [Hz] & \multicolumn{2}{c|}{$\sigma_j=\SI{0.1}{\micro s}$} &
\multicolumn{2}{c|}{\SI{1}{\micro s}} & \multicolumn{2}{c}{\SI{4}{\micro s}}\\
 & meas. & theory & meas. & theory & meas. & theory\\\midrule
100   & 84.03 & 84.04 & 64.07 & 64.04 & 52.02 & 52.00\\
1000  & 64.04 & 64.04 & 44.03 & 44.04 & 31.99 & 32.00\\
6500  & 47.76 & 47.78 & 27.79 & 27.78 & 15.75 & 15.74\\
10500 & 43.65 & 43.61 & 23.63 & 23.61 & 11.60 & 11.57\\
20000 & 38.03 & 38.02 & 18.01 & 18.02 &  6.20 &  5.97\\\bottomrule
\end{tabular}
\end{center}
The SNR drops by \SI{20}{dB} per decade of frequency and per decade of
jitter. Only at \SI{20}{kHz} with \SI{4}{\micro s} ($2\pi f\sigma_j=0.5$) does the
first-order approximation $e\approx\dot x\,\delta$ become inaccurate.

d) 12 bit + jitter at \SI{6.5}{kHz}: $\sigma_j=\SI{0.1}{\micro s}$:
\SI{47.8}{dB} (ENOB 7.6); \SI{1}{\micro s}: \SI{27.8}{dB} (ENOB 4.3);
\SI{4}{\micro s}: \SI{15.8}{dB} (ENOB 2.3). At \SI{100}{Hz} and
\SI{0.1}{\micro s} the full 12 bits remain (ENOB 11.97). Equality of the two
noise sources: $2\pi f\sigma_j=10^{-74/20}\Rightarrow f=\SI{31.8}{Hz}$ for
$\sigma_j=\SI{1}{\micro s}$. Conclusion: for kHz signals the sampling clock,
not the number of bits, decides the effective resolution; ECUs therefore
start knock conversions by hardware timers or sample continuously with a
fixed ADC clock.
\end{solution}

%-------------------------------------------------------------------------
\begin{exercise}{Time domain vs.\ crank-angle domain}{\JSLinux}
In \codefile{jslinux/sheet04/angle\_resample.c} the virtual engine
accelerates from \SI{1500}{rpm} to \SI{4500}{rpm} within \SI{1}{s}
(\texttt{e.rpm} is changed during the simulation). The ion-current signal is
sampled at $f_s=\SI{50}{kHz}$; the rising crank edges are time-stamped with a
resolution of \SI{2}{\micro s} (``input capture'' as in a real ECU), and the
cam level is latched at every edge.
\begin{tasks}
  \item Assign a crank angle to every edge: detect the gap (edge interval
    $>2\times$ the previous one), give the first edge after the first gap the
    angle \SI{0}{\degree} or \SI{360}{\degree} (cam level!), then add
    \SI{6}{\degree} per edge, \SI{18}{\degree} across the gap.
  \item Resample the signal onto the grid $\theta_m=\theta_0+m\cdot\SI{0.5}{\degree}$:
    for every $\theta_m$ interpolate the time $t(\theta_m)$ linearly between
    the two neighbouring edges, then the signal value $x(t(\theta_m))$
    linearly between two time samples. The program compares the grid with the
    true model angle: report the maximum and rms angle error and explain
    its size. Rebuild with \texttt{-DDT\_CAP=20e-6} (edges only known at the
    sample instants) and compare.
  \item The program correlates every engine cycle with the first one,
    (i) in the angle domain (1440 grid points per cycle) and (ii) in the time
    domain (segments with the sample count of the first cycle, each starting
    exactly at its cycle start). Interpret the table and the two plots.
  \item The knock-sensor signal contains crank-synchronous parts (combustion
    vibration $0\dots60^\circ$ after TDC) and fixed-frequency parts (valve
    impacts: \SI{8}{kHz}, $\tau=\SI{0.4}{ms}$). How do the two look in the
    angle domain during the ramp? Compute the angular length of $3\tau$ of a
    valve burst at \SI{1500}{rpm} and \SI{4500}{rpm}.
  \item Which angle step $\Delta\theta$ would be needed to sample the
    \SI{10.5}{kHz} knock mode directly in the angle domain at \SI{600}{rpm}
    and at \SI{6000}{rpm}? Which signals should an ECU therefore sample
    angle-based, which time-based?
\end{tasks}
\end{exercise}

\begin{solution}
a/b) Key part of the solution:
\inputcode[firstline=104,lastline=138]{solutions/jslinux/sheet04/angle_resample.c}
Output: 2899 crank edges, first reference edge \#57 at $t=\SI{38.52}{ms}$
(angle \SI{360}{\degree}, cam LOW), 35245 grid points. Angle error of the grid:
max \SI{0.054}{\degree}, rms \SI{0.022}{\degree}. The error is caused by the
time stamp resolution: an edge is registered up to $\Delta t_\mathrm{cap}$
late, i.e.\ $\omega\,\Delta t_\mathrm{cap}=27000\,\si{\degree/s}\cdot\SI{2}{\micro s}=\SI{0.054}{\degree}$
at \SI{4500}{rpm} -- exactly the maximum found. The linear interpolation
between edges (constant speed over \SI{6}{\degree}) is not the limit: the 1\,\%
second-order speed ripple changes $\omega$ by only $\approx0.2\,\%$ over a
tooth. With \texttt{-DDT\_CAP=20e-6}: max \SI{0.534}{\degree}, rms
\SI{0.222}{\degree} (and \SI{0.027}{\degree}/\SI{0.011}{\degree} for
\SI{1}{\micro s}): the error scales linearly with the capture resolution.
This is why ECUs time-stamp crank edges with a hardware capture unit instead
of the ADC sample clock.

c) Output (excerpt):
\begin{center}\small
\begin{tabular}{@{}rrrrr@{}}\toprule
cycle & rpm (mean) & samples/cycle & corr.\ time domain & corr.\ angle domain\\\midrule
1 & 1720 & 3488 & 1.000 & 1.0000\\
2 & 1919 & 3127 & $-0.030$ & 0.9986\\
5 & 2417 & 2482 & $-0.131$ & 0.9986\\
13 & 3407 & 1761 & $-0.035$ & 0.9985\\
22 & 4253 & 1411 & 0.043 & 0.9984\\\bottomrule
\end{tabular}
\end{center}
In the angle domain every cycle is (up to the \SI{10}{mV} noise of the model)
identical to the first one: correlation 0.998 for all 24 cycles. In the time
domain the correlation collapses already in the second cycle, although the
segments start exactly at the cycle start and the speed changed by only
12\,\%: the ion peaks (width a few degrees) at $120^\circ$, $300^\circ$, \dots{}
arrive \SIrange{1}{3}{ms} earlier and are narrower. The time plot shows the
last cycle compressed to 40\,\% (2.5 cycles in the window of the first
one), the angle plot shows two congruent curves. Crank-synchronous signals
are \emph{stationary in the angle domain} and non-stationary in the time
domain during speed changes; angle-based processing (synchronous averaging,
order analysis, Sheet~6) requires the resampling.

d) The combustion vibration is defined by the angle and stays congruent. A
valve burst is a resonance with fixed frequency and decay time: in the angle
domain its oscillation period is $\omega/f$ and changes with speed. $3\tau=\SI{1.2}{ms}$
corresponds to $9000\,\si{\degree/s}\cdot\SI{1.2}{ms}=\SI{10.8}{\degree}$ at
\SI{1500}{rpm} and \SI{32.4}{\degree} at \SI{4500}{rpm} (one \SI{8}{kHz} period:
\SI{1.1}{\degree} resp.\ \SI{3.4}{\degree}). Fixed-frequency resonances are
stationary in the time domain, crank-synchronous ones in the angle domain --
no single domain makes everything stationary.

e) The equivalent sampling rate is $f_\mathrm{eq}=\omega/\Delta\theta$;
$f_\mathrm{eq}>2\cdot\SI{10.5}{kHz}$ requires $\Delta\theta<\omega/\SI{21}{kHz}$:
\SI{0.17}{\degree} at \SI{600}{rpm} ($\omega=\SI{3600}{\degree/s}$) and
\SI{1.7}{\degree} at \SI{6000}{rpm}. A fixed angle step that is fine enough
at idle massively oversamples at high speed, a coarse one aliases at idle.
Hence: knock sensor (fixed resonances) is sampled time-based with fixed
$f_s$ and evaluated in a crank-angle \emph{window}; crank-synchronous
quantities (cylinder pressure, ion current, engine speed per tooth, misfire
detection) are sampled or resampled angle-based.
\end{solution}

\begin{deliverables}
4.1: derivations and numbers. 4.2: completed sketch, logic-analyzer screenshot
with sampling period and ISR duration, quiet-window results for three noise
levels. 4.3: SNR tables (quantisation, jitter, ENOB) with short discussion.
4.4: completed program, angle-error values for two capture resolutions, the
correlation table and the two plots with your interpretation.
\end{deliverables}
